% TOP_PROBLEM: {"id": 3063, "title": "Rota's unimodal conjecture", "category_id": 2, "category": {"id": 2, "name": "combinatorics", "display_name": "Combinatorics", "description": "Counting problems, graph theory, discrete structures.", "slug": "combinatorics", "order_index": 2, "created_at": "2026-07-31T15:26:25.671Z"}, "status": "open", "problem_number": "OPG-361", "set_id": 12}
% FIRST_POSTED: 2026-10-01
% ATTEMPT_STATUS: unresolved
% UnsolvedMath ID 3063; source catalogue code OPG-361
% !TeX program = lualatex
% GPT-6 Astra Ultra research attempt; independent partial work, 2026-10-01.
\documentclass[11pt,a4paper]{article}
\usepackage[margin=25mm]{geometry}
\usepackage{fontspec}
\setmainfont{lmroman10-regular.otf}[BoldFont=lmroman10-bold.otf,
 ItalicFont=lmroman10-italic.otf,BoldItalicFont=lmroman10-bolditalic.otf]
\usepackage{amsmath,amssymb,mathtools}
\usepackage{unicode-math}
\setmathfont{latinmodern-math.otf}
\AtBeginDocument{\let\setminus\smallsetminus}
\usepackage{xurl,hyperref}
\hypersetup{colorlinks=true,urlcolor=blue}
\setcounter{secnumdepth}{0}
\setlength{\parindent}{0pt}
\setlength{\parskip}{5pt}
\setlength{\emergencystretch}{3em}
\begin{document}
\section*{Rota's unimodal conjecture}\label{3063}
{\small UnsolvedMath ID 3063 · OPG-361 · Combinatorics · OpenGarden}
\subsection{Definitions and mathematical statement}
A finite matroid may be specified by a finite set $E$ and a rank function
$r:2^E\to\mathbb Z_{\ge0}$ satisfying $0\le r(X)\le|X|$,
monotonicity $X\subseteq Y\Rightarrow r(X)\le r(Y)$, and submodularity
$r(X\cup Y)+r(X\cap Y)\le r(X)+r(Y)$. Put $R=r(E)$.
A flat, or closed set, is a subset $F\subseteq E$ such that
$r(F\cup\{e\})>r(F)$ for every $e\in E\setminus F$.
For $0\le i\le R$, let $w_i$ be the number of flats of rank $i$.

The first conjecture asserts that this sequence is unimodal: there is
an index $j\in\{0,\ldots,R\}$ such that
\[
 w_0\le\cdots\le w_j\ge w_{j+1}\ge\cdots\ge w_R.
\]
The accompanying stronger conjecture is log-concavity:
\[
 w_i^2\ge w_{i-1}w_{i+1}\qquad(1\le i<R).
\]
Both assertions quantify over every finite matroid, without any assumption
of representability over a field. For rank zero or one, the displayed
log-concavity conditions are empty.

The numbers count flats themselves, each once, even if the matroid has
loops or parallel elements. They are the Whitney numbers of the second
kind of the lattice of flats. They are not the numbers of independent
sets of each cardinality, and are not the absolute coefficients of the
characteristic polynomial. Results concerning those different sequences
do not automatically settle either assertion here.
\subsection{Short English statement}
Are the numbers of flats of successive ranks in every finite matroid unimodal, and more strongly log-concave?
\subsection{Research attempt: a low-rank incidence certificate and a status correction}
\paragraph{Literature check and honest target status.}
GPT-6 Astra Ultra, 1 October 2026. The canonical definition above records
the historical conjectures and is preserved verbatim. The July 2026
preprint of Divoux, Lowen and Wang [S3] gives counterexamples to
unimodality, building on counterexamples to the stronger log-concavity
assertion. Thus the catalogue's historical ``open'' label must not be
read as a fresh statement that the general conjecture has no disproof.
This attempt does not independently certify their whole construction.
It gives a self-contained rank-three verification and identifies why
it cannot support an induction proving the historical conjecture.

\paragraph{1. Simplification does not change the flat counts.}
Delete loops and replace each parallel class by a single representative.
Every flat of the original matroid contains every loop and, whenever
it contains a nonloop, contains the entire parallel class of that
element. Conversely the inverse image of a flat in the simplified
matroid, with loops adjoined, is a flat of the same rank. This is a
rank-preserving bijection of flat lattices. It therefore suffices to
work with simple matroids. Ranks zero, one and two have flat-count
sequences $(1)$, $(1,1)$, and $(1,p,1)$ respectively, where $p\ge2$
in the last case, so satisfy both inequalities.

\paragraph{2. Rank three: two incidence counts.}
Let there be $p$ points (rank-one flats) and $l$ lines (rank-two
flats). Each pair of different points spans a unique line, and every
line contains at least two points. Distinct lines cannot contain the
same pair. Consequently
\[
 l\le\binom p2<p^2.
\]
Every point lies on at least two lines. Indeed, choose a second point
and then a point outside their rank-two closure, possible because the
matroid has rank three. The lines through the original point and these
two other points are distinct.

Let $B$ be the $p\times l$ point-line incidence matrix, considered over
the real numbers even if the matroid has no representation over a field.
If $d_x$ counts lines through point $x$, unique pair incidence gives
\[
 BB^{\mathsf T}=J+\operatorname{diag}(d_x-1),
\]
where $J$ is the all-one matrix. For every nonzero real vector $z$,
\[
 z^{\mathsf T}BB^{\mathsf T}z
 =\left(\sum_xz_x\right)^2+\sum_x(d_x-1)z_x^2>0.
\]
The strict inequality uses $d_x\ge2$. Thus $BB^{\mathsf T}$ has
rank $p$, so $l\ge p$. The full sequence is $(1,p,l,1)$;
$p^2\ge l$ and $l^2\ge p$ prove log-concavity and hence unimodality.
This elementary incidence argument is classical in character, and no
novelty is claimed.

\paragraph{3. An attempted higher-rank generalization fails precisely.}
Try to replace $B$ by the incidence matrix between rank-$i$ and
rank-$(i+1)$ flats. At rank one in a rank-three matroid, every pair of
points lies on exactly one line, which made every off-diagonal Gram
entry equal to one. For $i\ge2$, two rank-$i$ flats can have join
of rank greater than $i+1$, and then have no common rank-$(i+1)$
upper flat. Others can have a common upper flat. There is no resulting
formula $BB^{\mathsf T}=J+D$ with $D$ positive diagonal. Moreover,
injectivity of one incidence map would give only an inequality between
two flat counts, not the three-term log-concavity inequality.

\paragraph{Verification and remaining task.}
The proof needs neither representability nor an assumption on the
number of elements in a line. For the Fano rank-three matroid it gives
$(1,7,7,1)$; for the uniform rank-three matroid on $p$ elements it
gives $(1,p,\binom p2,1)$, consistent with both inequalities. The
historical universal assertion cannot be concluded from these small
examples. Independently reconstructing the cited high-rank counterexample,
including its matroid axioms and every flat count, is outside what this
attempt has proved; this is an audit gap, not evidence that the
historical conjecture should still be pursued as true.

\subsection{Final status}
UNRESOLVED as an independent full reconstruction of the catalogue
question; the rank-at-most-three subcase is proved. The general
historical conjecture is reported disproved in [S3], which is credited
explicitly. No new disproof or full verification of [S3] is claimed.

\subsection{Sources}
\begin{itemize}
\item[S1] ULAM.AI and the UnsolvedMath Contributors, supplied problems.json, record 3063 (OPG-361), OpenGarden collection. Catalogue identity and source transcription; CC BY 4.0.
\url{https://huggingface.co/datasets/ulamai/UnsolvedMath}.
\item[S2] Open Problem Garden, Rota's unimodal conjecture. Original problem and discussion; the mathematical formulation below is expanded from the supplied source record.
\url{http://www.openproblemgarden.org/op/rotas_unimodal_conjecture}.
\item[S3] A. Divoux, C. Lowen and S. Wang, \emph{Matroid flat counts are not unimodal}, arXiv:2607.22515v1, 24 July 2026. Abstract and paper checked 1 October 2026; source of the status qualification.
\url{https://arxiv.org/html/2607.22515v1}.
\end{itemize}
\end{document}
